\myaddRC{\section{Preliminaries}\label{sec:preliminaries}
}
% In this section, we introduce the graph notation and the definitions used to compute per‑clique $(r,s)$‑core numbers; hierarchy/connected components are deferred to later sections.

%--------------------------------------------------
% 符号表
%--------------------------------------------------
% \begin{table}[ht]
%   \centering
%   \caption{Notation used throughout the paper}\label{tab:notation}
%   \begin{tabular}{ll}
%     \hline
%     \textbf{Symbol} & \textbf{Meaning}\\
%     \hline
%     $G=(V,E)$     & simple, unweighted graph \\
%     $G[V']$       & subgraph of $G$ induced by $V' \subseteq V$\\
%     $n = |V|$     & \#vertices of $G$\\
%     $m = |E|$     & \#edges of $G$\\
%     $\deg(v)$     & degree of vertex $v$\\
%     $C\subseteq V$ & a clique (context specifies its size)\\
%     $\deg_{(s)}(R)$ & $s$-clique degree of an $r$-clique $R$ in $G$\\
%     $G_k$         & a $k$-\((r,s)\)-nucleus (defined below)\\
%     $\kappa_{s}(R)$ & \((r,s)\)-nucleus core number of $r$-clique $R$\\
%     \hline
%   \end{tabular}
% \end{table}

%--------------------------------------------------
% 基本概念
%--------------------------------------------------
\stitle{Graph Notations.} 
We focus on a \emph{simple}, \emph{unweighted} graph, defined as a pair $G=(V,E)$ where $V$ is the set of vertices and $E\subseteq V\times V$ is the set of undirected edges. 
%
We use $N_G(v)=\{u\in V \mid (u,v)\in E\}$ to denote the set of neighbors of a vertex $v\in V$ in $G$, and $\deg_G(v)=|N_G(v)|$ to denote its degree.
% 
We use $G[V']$ to denote the subgraph \emph{induced} by a vertex subset $V'\subseteq V$. 
%
We store a graph $G$ in the classical Compressed Sparse Row (CSR) format, consisting of two integer arrays: \texttt{offsets} of length $|V|{+}1$ and \texttt{edges} of length $|E|$.

\begin{definition}[$k$-Clique]\label{def:k-clique}
A \emph{$k$-clique} is a vertex set $C\subseteq V$ with $|C|=k$ such that $(u,v)\in E$ for every $u,v\in C$. 
\end{definition}

% \begin{example}
%     Each vertex is a $1$-clique, each edge is a $2$-clique, and each triangle is a $3$-clique.
% \end{example}

\begin{definition}[$(r,s)$-Clique Degree]\label{def:r-s-degree}
Given two integers $0 < r < s$, 
the \emph{$(r,s)$-clique degree} of an $r$-clique $R$ in $G$ is defined as
\[
  \deg_{(s)}^{G}(R)
  \;=\;
  \bigl|\,\{\,S \subseteq V \mid R \subseteq S, \; S \text{ is a } s\text{-clique in } G\,\}\,\bigr|.
\]
\end{definition}
% 说一下这个也叫做 $s$-clique support
We also refer to $\deg_{(s)}^{G}(R)$ as the \emph{$s$-clique support} of $R$ in $G$.
When the graph $G$ is clear from the context, we simply write $\deg_{(s)}(R)$.

%--------------------------------------------------
% nucleus 核心定义
%--------------------------------------------------

% \begin{definition}[{$k$-admissible $s$-clique}]
% Given integers $0<r<s$ and $k>0$.  
% An induced $s$-clique $S$ in $G$ is called \emph{$k$-admissible}
% if every $r$-clique $R\subseteq S$ satisfies $\deg_{(s)}(R)\ge k$.
% \end{definition}

\begin{definition}[{$k$-admissible $s$-clique}]
Given integers $0 < r < s$ and $k > 0$,  
an induced $s$-clique $S$ in $G$ is called \emph{$k$-admissible} 
if every $r$-clique $R \subseteq S$ satisfies $\deg_{(s)}(R) \ge k$.
\end{definition}


% \begin{definition}[{$r$-clique-support}]
% The \emph{$r$-clique-support} of an $s$-clique $S$ in $G$ is defined as
% \end{definition}

\begin{definition}[$k$-$(r,s)$-nucleus]\label{def:k-rs-nucleus-new}
Let $\mathcal S_k$ be the set of all $k$-admissible $s$-cliques of $G$.
For any subset $\mathcal S\subseteq\mathcal S_k$, denote
$
   G[\mathcal S] \;=\; G\, \!\Bigl[\,
        \bigcup_{S\in\mathcal S} S
   \Bigr].
$
A subgraph $G_k = G[\mathcal S]$ is a \emph{$k$-$(r,s)$-nucleus} if
\begin{enumerate}[leftmargin=*]
  \item $G_k$ is \emph{$s$-connected}: for any two $r$-cliques
        $R,R'\subseteq V(G_k)$, there is a sequence
        $R=R_1,\dots,R_t=R'$ with
        $R_i\cup R_{i+1}\subseteq S_i$ for some
        $S_i\in\mathcal S$;
  \item $G_k$ is \emph{maximal}, i.e., no proper superset
        $\mathcal{S}' \supset \mathcal{S}$ still satisfies condition~(1).
\end{enumerate}
\end{definition}

% \stitle{Intuition.}
A $k$-$(r,s)$-nucleus is a maximal subgraph in which 
every $r$-clique participates in at least $k$ distinct $s$-cliques, 
and all such $r$-cliques are tightly connected through shared $s$-cliques. 
% The formal definition is given below.

\begin{definition}[Core Number]\label{def:r-clique-core-number}
Given two integers $0 < r < s$,  
the \emph{$(r,s)$-nucleus core number} (or \textit{coreness}) of an $r$-clique $R$, denoted by $\kappa_s(R)$, 
is the largest integer $k$ such that there exists a $k$-$(r,s)$-nucleus $G_k$ in $G$ containing $R$. 
Equivalently, it is the largest integer $k$ for which $R$ belongs to some $k$-$(r,s)$-nucleus $G_k$ in $G$.
\end{definition}

%--------------------------------------------------
% s-clique connectivity
%--------------------------------------------------
% \begin{definition}[$s$-Connectivity]\label{def:s-connect}
% Fix integers $0<r<s$.  
% Two $r$-cliques $R$ and $R'$ in a graph (or subgraph) $H$ are \emph{$s$-connected}
% if there exists a sequence of $r$-cliques
% $
%   R=R_1, R_2, \dots, R_t=R'
% $
% such that for each $i\in\{1,\dots,t-1\}$ there is an induced $s$-clique
% $S_i\subseteq V(H)$ with $R_i\cup R_{i+1}\subseteq S_i$.
% The sequence $\langle R_1,\dots,R_t\rangle$ is called an \emph{$s$-clique path}.
% \end{definition}

% \begin{definition}[$s$-Connected Component]\label{def:s-component}
% For fixed $0<r<s$, the \emph{$s$-connected component} of an $r$-clique
% $R$ (with respect to $s$) is the vertex-induced subgraph
% \[
%   \mathrm{Comp}_s(R) \;=\;
%   G\!\Bigl[\,
%     \bigcup_{R'\,:\;R'\text{ is }s\text{-connected to }R} R'
%   \Bigr].
% \]
% Equivalently, it is the maximal subgraph $H$ containing $R$
% such that any two $r$-cliques in $H$ are $s$-connected within $H$.
% $s$-Connectedness is an equivalence relation on the set of $r$-cliques; hence
% $s$-connected components form a partition of all $r$-cliques in $G$.
% \end{definition}

\sstitle{Problem statement.}
Given a graph $G$ and integers $0 < r < s$, the goal is to compute the $(r,s)$-nucleus core number $\kappa_s(R)$ for every $r$-clique $R$ in $G$. (The hierarchy construction of $k$-$(r,s)$ nuclei will be formally discussed in Section~\ref{sec:hierarchy}.)


